\section{Effective comparison of entry maps}\label{sec:entry-comparison}

Residues alone do not decide inequalities between sums of collapse values.
The order argument also needs exact comparisons and collision conditions
for entry norms. This section supplies them for the chosen controls.
All ordinal additions in interval endpoints and relative states are ordinary
ordinal additions. Natural sum and product remain explicitly marked.

For a proper interval \(P=(\lambda,\mu)\), let \(\delta_P\) be the
unique nonzero limit with \(\mu+\delta_P=\lambda\). Put
\(E_P=E_{\lambda,\mu}\) and \(\tau_P=\tau_{\lambda,\mu}\).
An effective sufficient domain for entry is
\begin{equation}\label{eq:affordable-tail}
 t\ge\max\{2,N(\lambda),N(\mu)\}.
\end{equation}
Indeed, before a hit the target is an affordable predecessor, so the
descent cannot skip below it; strict ordinal descent forces a hit.
This argument also covers \(\mu=0\). A smaller genuine input may
skip a positive target and is not assigned an invented entry value.

\subsection{Relative descent and the two comparison rules}

For \(y>0\), let \(e(y)\) be its leading Cantor exponent. Define
\[
 s_\mu(e)=N(\mu_{\ge e}),\qquad s_\mu(e(0))=N(\mu).
\]
The subscript means restriction to Cantor terms of exponent at least
\(e\); the convention at zero is separate. Directly from Cantor
normal form,
\begin{equation}\label{eq:relative-norm}
 N(\mu+y)=N(y)+s_\mu(e(y)),\qquad 0\le s_\mu(e(y))\le N(\mu).
\end{equation}
The terms of \(\mu\) below the leading exponent of \(y\) are absorbed;
the contributions at equal exponents add. Before entry, write the actual
state uniquely as \(\mu+y_i\). With \(b_i=f^i(t)\), its exact recursion is
\begin{equation}\label{eq:relative-descent}
 y_0=\delta_P+t-N(\lambda),\qquad
 y_{i+1}=\max\{y<\operatorname{lim}(y_i):
          N(y)+s_\mu(e(y))\le b_{i+1}\}.
\end{equation}
Here \(\operatorname{lim}\) removes the finite tail. The first finite
\(y_i\) is the entry event, and \(N(y_i)\le b_i\).

Consider two recursions of this form with profiles \(s,s'\) bounded by
\(S\), budgets \(b,b'\), and limit parts \(U_i,V_i\).
Two distinct budget comparisons are needed.

\begin{lemma}[Propagation]\label{lem:propagation}
If \(V_{i+k}>U_i=\rho>0\) and
\(b'_{i+k+1}\ge b_i+S\), then
\(V_{i+k+1}\ge\rho>U_{i+1}\).
If instead \(V_{i+k}=U_i\) and
\(b'_{i+k+1}\ge b_{i+1}+S\), then
\(V_{i+k+1}\ge U_{i+1}\).
For equal next budgets the latter conclusion also follows from
\(s'\le s\) on all eligible candidates.
\end{lemma}
\begin{proof}
In the strict case, \(\rho\) is below the upper limit part and
\(N(\rho)+s'(e(\rho))\le N(y_i)+S\le b_i+S\).
It is therefore an upper candidate. The lower next state is strictly
below \(\rho\). In the equal-limit case every lower candidate costs
at most \(b_{i+1}+S\) under the upper profile. Candidate inclusion
also proves the equal-budget alternative.
\end{proof}

If the strict rule persists until the lower process stops at \(n\), the
upper process stops no earlier than \(n+k+1\): at the last lower event
its upper limit part is still positive. The weak rule gives the bound
\(n+k\). Notice that the strict rule compares the next upper budget
with the current lower budget. Using the equal-limit budget condition
for two identical budget orbits would lose the required slack.

\begin{lemma}[Gap and input comparison]\label{lem:gap-input}
For two fixed intervals \(P,Q\), put
\(S=\max\{1,N(\mu_P),N(\mu_Q)\}\). On effective tails:
\begin{enumerate}[label=(\roman*)]
\item If \(\delta_P>\delta_Q\), then
\(\tau_P(t)>\tau_Q(t)\) and \(E_P(t)>E_Q(t)\).
\item If \(\delta_P\ge\delta_Q\) and \(t>u\), then
\(\tau_P(t)\ge\tau_Q(u)\) and \(E_P(t)>E_Q(u)\).
\item For fixed \(k\ge1\), if
\(\delta_P\ge\delta_Q+\omega k\) and \(u+S\le f^k(t)\), then
\(\tau_P(t)\ge\tau_Q(u)+k\) and \(E_P(t)>E_Q(u)\).
\end{enumerate}
\end{lemma}
\begin{proof}
Take initially
\(T\ge\max\{2,S,N(\lambda_P),N(\lambda_Q)\}\).
The estimate \(f(v)\ge v+S\) for \(v\ge T\) propagates the strict
initial gap in (i). For (ii), adjacent growth gives
\(f(t)\ge f(u)+S\); monotonicity and noncontraction preserve this
gap at every later iterate. The weak rule propagates the initial weak
limit comparison. The larger starting input then makes the terminal
norm strictly larger.

For (iii) enlarge the tail by \(2k\). During its first \(k\) events
the upper process retains the successive reserves
\(\delta_Q+\omega(k-j)\). At event \(j<k\) the next reserve costs
at most \(N(\delta_Q)+2(k-j-1)+S\), below the next norm budget.
At lag \(k\), compare with the lower process using the weak rule and
\(f^k(t)\ge u+S\). This yields the count bound. Since \(S\ge1\),
the terminal inequality is strict. The positive unit is necessary when
both targets are zero: with \(P=(\omega\cdot2,0)\), \(Q=(\omega,0)\),
and \(u=f(t)\), the terminal norms coincide.
\end{proof}

Distinct limit gaps differ by at least an \(\omega\)-tail. Thus fixed
input translations do not change their eventual order, by (iii) with
\(k=1\). For equal gaps, unequal fixed input translations are ordered
by (ii). It remains to compare equal gaps at one common input.

\subsection{Equal gaps: a finite phase calculation}

Suppose \(\delta_P=\delta_Q=\delta\). Set
\[
 R=20+N(\delta)+N(\mu_P)+N(\mu_Q),\qquad
 M=\operatorname{lcm}(1,\ldots,R+1).
\]
Choose the input tail above \(100(R+2)^4\), \(M\), and \(B_f\),
and above the finitely many affordability bounds used below. In
particular every next tail budget exceeds the current prefix cost and
all fixed shifts by more than \((R+3)^3\).
All fixed exponents and breakpoints occurring in this argument have
norm at most \(R+1\). Write \(s=s_{\mu_P}\), \(s'=s_{\mu_Q}\),
and \(d=s'-s\).

Suppose two still equal limit parts have leading exponent \(E\), and
write a prefix-retaining step as
\(\rho=P_0+\omega^a\), with \(P_0\ne0\) retaining that leading
exponent. The exact next tail is
\begin{equation}\label{eq:local-phase}
 P_0+M(\omega^a,X),\qquad X=b_{i+1}-s(E)-N(P_0).
\end{equation}
The growth margin covers the retained prefix and a positive remainder;
any state below the prefix is inferior, and a tail at least
\(\omega^a\) is outside the interval. This proves the formula.

For \(a=1\) the new tail is finite, so a change of \(X\) does not
change the limit part. If \(a\) is a nonzero limit, the leading
exponent of the new tail is \(M(a,X-1)\), strictly increasing with
\(X\) by saturation. For \(a=a_0+1\), the exact recursion is
\begin{equation}\label{eq:successor-tail}
 M(\omega^a,X)=\omega^{a_0}q+M(\omega^{a_0},r),
 \quad q=\lfloor X/v\rfloor,\quad r=X\bmod v,
 \quad v=1+N(a_0).
\end{equation}
In particular, at \(a=2\) two different budgets have the same limit
part precisely when their difference is one and the larger budget is odd.

\begin{lemma}[A retained thick block forces separation]\label{lem:thick}
Suppose two undiverged relative descents process a common block
\(\omega^a q\), where \(a\ge2\) is fixed, \(N(a)\le R\), and
\(q>R+2\), while their whole-state leading exponent, and hence their
two distinct shifts \(c<c'\), remain fixed. Their limit parts diverge
before this block is processed.
\end{lemma}
\begin{proof}
If a lower remainder already creates a divergence, the conclusion holds.
Otherwise inspect the common next top visit
\(\rho=P_0+\omega^a(q-1)\). In each absolute descent the preceding
top visit lies at least \(\omega^2\) above this target. Both target
funnels are affordable. The long-entry residue formula, applied modulo
\(2M\) at the common preceding norm, gives
\[
 \left\lceil\frac{c+N(\rho)}2\right\rceil
 \equiv
 \left\lceil\frac{c'+N(\rho)}2\right\rceil\pmod M.
\]
The difference of these rounded quantities is smaller than \(M\),
so they are equal. It follows that \(c'=c+1\) and
\(c+N(\rho)\) is odd. At the next event the common budget \(u\)
satisfies, by the modular decrement,
\[
 u\equiv c+N(\rho)-1\pmod{2M}.
\]
Peel one unit, leaving \(P^-=P_0+\omega^a(q-2)\). The large
coefficient ensures that this prefix retains the original leading
exponent and hence the fixed shifts. The two local tail budgets are
\(X=u-c-N(P^-)\) and \(X-1\), with
\[
 X\equiv N(a)\pmod{2M}.
\]
If \(a\) is a limit, saturation makes the next leading exponents
different. If \(a=a_0+1\), then \(N(a)=1+N(a_0)=v\), and
\(v\mid M\). Thus \(v\mid X\). Formula
\eqref{eq:successor-tail} gives coefficients \(X/v\) and \(X/v-1\)
of the infinite \(\omega^{a_0}\)-block. The limit parts differ.
\end{proof}

This proof uses the moving residue of our control. A proof assuming that
every updated budget has one fixed residue would not apply to
\eqref{eq:factorial-control} or \eqref{eq:near-control}.

\begin{proposition}[An effective equal-gap certificate]\label{prop:equal-gap}
For two fixed equal-gap intervals one can compute a threshold and decide
whether their entry maps are identical on that tail or have one fixed
strict order there.
\end{proposition}
\begin{proof}
Here is the complete finite calculation. Let \(E=e(\delta)\).
During the initial phase, each nonleading term of \(\delta\) is
peeled one unit at a time, followed by all but one unit of its leading
term. The generated smaller blocks are processed before the next
original unit is peeled. Call a peeled unit sensitive if its exponent
is at least two.

If \(d(E)=0\), all prefix-retaining steps agree exactly. If
\(d(E)\ne0\), a sensitive limit exponent gives immediate divergence
by \eqref{eq:local-phase}. A successor exponent at least three either
gives immediate divergence or creates a common thick block, to which
Lemma~\ref{lem:thick} applies. Only exponent two can fail to separate.
Its coincidence condition is a shift difference of one and an odd larger
local budget. Consecutive original \(\omega^2\)-units have prefix
norms differing by three, whereas all updated budgets are even. Their
coincidence parities alternate, so two such units cannot both fail.
Thus the only nonempty sensitive initial phases which can fail to
separate have shapes
\[
 \delta=\omega^E+\omega^2+\omega c\quad(E\ge3),\qquad
 \delta=\omega^2\cdot2+\omega c,
\]
and their exact exception is
\[
 |d(E)|=1,\qquad \min\{s(E),s'(E)\}+N(P_0)\text{ odd},
\]
where \(P_0\) is the retained prefix at that one sensitive step.
If there is no sensitive step, the initial phase also fails to separate.
Every initial divergence has sign \(d(E)\) for \(E_P-E_Q\), by
candidate inclusion and Lemma~\ref{lem:propagation}.

If the initial phase has not separated, refine jointly the finitely many
breakpoint partitions of \(s\) and \(s'\). Search for the highest
piece on \(2\le e<E\) on which \(d(e)\ne0\). This is a finite
Cantor-normal-form calculation. An untruncated piece is right-closed,
so use its included right endpoint below \(E\). Only the top piece
clipped by the excluded endpoint \(E\) needs a different choice:
if \(E\) is a successor, use its predecessor when this lies in the piece;
if \(E\) is a limit and \(a\) is the lower edge, use the fixed exponent
\(a+1<E\). Discard an empty clipped piece and continue with the next one.

To justify that this finite piece determines the descent, note first
that phase exponents strictly decrease. A successor exponent drops to
its predecessor. A jump from a limit exponent cannot skip a fixed
\(\zeta\) of norm at most \(R+1\), since
\(\omega^{\zeta+1}\) is an affordable candidate. If the highest
nonzero piece has a maximum \(\beta\), the two descents agree until
that piece is reached. They enter it by a drop from \(\beta+1\):
a jump to the bounded \(\beta\) would have the larger affordable
candidate \(\omega^{\beta+1}\). The drop produces a common block
\(\omega^\beta q\) with \(q>R+2\), unless it has already diverged.
Lemma~\ref{lem:thick} then gives sign \(d(\beta)\).

A highest piece without a maximum is necessarily the top truncated
piece \((a,E)\), where \(E\) is a limit. On leaving the initial
pure \(\omega^E\)-state, a fixed exponent above \(a\) is affordable,
so both winning exponents lie within that piece. There the shifts are
constant, and the two exponents are
\(M(E,b_{i+1}-c-1)\) and \(M(E,b_{i+1}-c'-1)\).
Saturation separates them, with sign determined by the smaller shift.

If no such piece exists, all phases at exponents at least two finish
with identical limit parts and common budget index. For \(E\ge2\)
the final pure \(\omega^2\)-phase has next coefficients
\[
 b_{i+1}/2-\lceil s(1)/2\rceil,\qquad
 b_{i+1}/2-\lceil s'(1)/2\rceil.
\]
Their difference decides the order. Equal coefficients give simultaneous
termination. For \(E=1\), the original \(\omega\)-coefficient
already gives the same exact count in both descents.

The algorithm therefore uses, in order: the nonexceptional initial
phase; the highest nonzero profile piece on \([2,E)\); the difference
\(\lceil s'(1)/2\rceil-\lceil s(1)/2\rceil\); and otherwise
identity. All margins used above are covered by the common
tail. Equal stopping times give equal entry norms by
\eqref{eq:entry-orbit}; first limit-part divergence gives strict order
by propagation. This proves both completeness and effectiveness.
\end{proof}

\subsection{Cocycles, inverse thresholds, and collisions}

For fixed \(0\le\mu<\kappa<\lambda\), entry through \(\kappa\) requires
its affordability. On a tail satisfying all three endpoint norm bounds,
literal continuation gives
\begin{equation}\label{eq:cocycle}
 \begin{split}
 E_{\lambda,\mu}(t)&=
 E_{\kappa,\mu}(E_{\lambda,\kappa}(t)),\\
 \tau_{\lambda,\mu}(t)&=\tau_{\lambda,\kappa}(t)+
 \tau_{\kappa,\mu}(E_{\lambda,\kappa}(t)).
 \end{split}
\end{equation}
Genuine definedness of an exit alone is insufficient: at \(\lambda=\omega^2\),
\(t=3\), and \(B=f(3)\), the first predecessor is
\(\omega\cdot(B/2)\), which skips the fixed target
\(\kappa=\omega\cdot(B/2+1)\). A uniform sufficient guard for the
cocycle is \(t\ge\max\{2,N(\lambda),N(\kappa),N(\mu)\}\).

\begin{proposition}[Exact inverse-threshold table]\label{prop:threshold}
Given proper intervals \(P,Q\), comparisons of \(E_P(x)\) and
\(E_Q(y)\) have effective finite descriptions using comparisons of raw
inputs with fixed entry maps on the other input. Equality cases are
decided by the same description.
\end{proposition}
\begin{proof}
For equal gaps, Lemma~\ref{lem:gap-input} handles unequal inputs and
Proposition~\ref{prop:equal-gap} handles equal inputs. Suppose
\(\delta_Q<\delta_P\). Put
\[
 \kappa=\mu_P+\delta_Q,\quad
 \widetilde Q=(\kappa,\mu_P),\quad
 H=(\lambda_P,\kappa),\quad h=E_H(x).
\]
All these endpoints are effective fixed Cantor normal forms, and
\(\mu_P<\kappa<\lambda_P\). Impose their norm bounds and the relevant
comparison thresholds. The cocycle gives
\(E_P(x)=E_{\widetilde Q}(h)\).
Let \(\sigma\) be the eventual sign of
\(E_{\widetilde Q}-E_Q\) at a common input, supplied by the equal-gap
certificate. Input comparison brackets
\[
 E_Q(h-1)<E_{\widetilde Q}(h)<E_Q(h+1).
\]
Enlarge the tail so \(h-1\) is within all required domains. The exact
comparison of \(E_Q(y)\) with \(E_P(x)\) is consequently
\begin{center}\small
\begin{tabular}{@{}llll@{}}\toprule
\(\sigma\)& smaller & equal & larger\\\midrule
\(0\)&\(y<h\)&\(y=h\)&\(y>h\)\\
\(1\)&\(y\le h\)&never&\(y\ge h+1\)\\
\(-1\)&\(y\le h-1\)&never&\(y\ge h\)\\\bottomrule
\end{tabular}
\end{center}
Swap roles for the other gap order. If both original targets are zero,
\(\widetilde Q=Q\), so the middle certificate is identity.
Fixed input translations remain explicit in this table.

Every proper entry value belongs to the range of \(f\). Two distinct
such values \(U<V\) satisfy \(V\ge2^{U+1}\), by adjacent growth.
On a larger effective tail any fixed output offsets are therefore
irrelevant except at a collision, where their comparison is an ordinary
integer test. This also supplies the asserted descriptions with fixed
output offsets.
\end{proof}

The construction introduces the interval \(H\) and also the certificate
partner \(\widetilde Q\). Both must be retained when computing the
finite collection of source norms and thresholds. No inverse function
is added as a new primitive.
